% !TeX root = ../papers/nc-p3.tex
\section{二本の直線を通る分岐と純関係のない分岐の三次曲線}
\label{sec:two-line-branch-blocks}
\begin{proposition}\label{prop:two-line-branch-block}
曲面射影のない標準正則窓で、分岐曲線列が
$p,x_1,x_2,z$ と二つの動く位置を持つとする。
二つの直線に制限した四テンソルは $[x,y]h$ であり、$h\in\ker z\otimes\ker p$ は階数3である。
$p,z$ の中央支持次元を $a,b$ とし、$r=\operatorname{rank}w(p,z)$ とすると
\[
                        \rho\ge 6+a+b-r,\qquad r\le2 .
\]
特に純関係がなければ $\rho\ge12$ である。
$\rho=8$ の場合には $a=b=r=2$ となり、二つの直線の消滅二空間どうしに
二項対称二乗と同型な三次元の純関係の線型包がある。
\end{proposition}
\begin{proof}
両直線のパラメータを移動で揃える。$R_1$ が対角で消えるので制限表示が得られ、
両定点との隣接消滅から $h$ の支持は表示通りである。
中間辺の一般直交補はここでは全点で相補的である。実際 $x_1$ の逆支持は常に $\ker p$、
$x_2$ の順支持は常に $\ker z$ である。従って $R_2\simeq R_0^*$ の六空間を使う
定理\ref{thm:finite-branch-corridors}の直交束の列は
\[
          0\longrightarrow3\mathcal O(-1)\longrightarrow6\mathcal O
           \longrightarrow3\mathcal O(1)\longrightarrow0
\]
となる。左右を双対にしても同様であり、$H^0,H^1$ で $\mathcal O(-1)$ が消えるので
$R_0$ の位置1の直線への制限も、$R_2$ の位置2の直線への制限も六次元のままである。
これらは互いに異なる側の切断なので、同時に制限しても隣接切断の階数6は保たれる。
$[x,y]h$ のその階数は $2\operatorname{rank}h$ であり結論を得る。

反対切断でもこの表示は六次の可逆部分行列を与える。
位置0は $kp$ の補空間、位置2は直線上の二点座標の包を使う六行であり、
位置1の二点座標と位置3の $kz$ の補空間が六列である。
追加で $p\otimes V_2^*$ の四行、$V_1^*\otimes z$ の四列を取る。
追加行は先の六列で零（最初の直線は $p$ の後続）、追加列も先の六行で零である。
従って可逆部分を消去しても追加行・列の全階数はそれぞれ $a,b$、その交叉部分行列は $w(p,z)$ のままである。
一般の行列の階数は指定行群の階数と指定列群の階数の和から両群の交叉部分の階数を引いた値以上である。
これは指定行群から独立行を取り、残りの行の指定列への射影を数えれば従う。
$r\le2$ は両端の直線支持である。
純関係がなければ両側から系\ref{cor:pure-free-branch-rank-nine}を使える。
$\rho=8$ なら中央支持は少なくとも2なので全ての等号が必要である。
$p$ の $2\times2\times3$ 縮約の三次元関係像は $L_2$ の二項対称二乗であり、
$r=2$ によってその中央二空間は第二直線の消滅二空間に一致する。
\end{proof}

\begin{proposition}\label{prop:pure-free-branch-cubic}
純関係のない曲面射影のない正則窓で、任意の順側中心 $p$ の直後の辺の点対は、
第一点がその後続直線にあるか、第二点が $p$ から一意に決まる有理正規三次曲線にある。
一位置の順側中心は高々二点で、二点あればその位置の各後続直線の次の曲線は
相手に付随する三次曲線であり、次位置に順側中心はない。
反対向きも同様である。曲線列の最初の二曲線がどちらも直線なら、動く位置はその二つで尽きる。
\end{proposition}
\begin{proof}
系\ref{cor:pure-free-branch-rank-nine}の $L_3$ 表示で、位置2の文字は二項三次となり、
その点 $z$ との縮約は二行三列のHankel行列である。
その階数1以下軌跡はVeronese三次曲線である。外では行支持が $p$ の順支持全体を張るから、
$z$ の全前点は $p$ の後続直線上にある。
次位置に別の順側中心があれば、その後続直線はこの整三次曲線に入らないので
その中心自体が $p$ の後続直線に属する。

同位置の相異なる二中心の直線は相異なる。一方の直線の次の像は一定ではなく、
上の条件より相手の三次曲線に一致する。その二つの三次曲線も相異なる。
同じなら一般点が二つの異なる前点を持つか、元の相異なる二直線に同じ一般前点が
同時に属することになって不可能である。第三の中心があれば第二と第三の三次曲線が
第一の直線の次像として一致して矛盾する。二中心のとき各直線上で $d_2=3$ ゆえ
$f_1=0$ であり、次位置の中心はそこにない。

$d_2=1$ なら $f_1=2$ なので次位置にも順側中心がある。$m\ge3$ とすると、
$x_2$ の像はその新中心の後続直線ではない（一般前点は動く）。
よって $x_3$ はその三次曲線となる。一方 $f_2=3d_2-3-d_3$ は負になってしまう。
これで $m=2$ である。
\end{proof}

\begin{corollary}\label{cor:two-line-secant-roots}
純関係のない二直線列の両側で、最初の直線は終中心 $z$ の反対三次曲線の割線または接線であり、
その長さ二の交叉が最初の直線の順側分岐因子 $f_1$ をスキームとして与える。
最後の直線でも $p$ の三次曲線について逆向きに同じである。
\end{corollary}
\begin{proof}
直線の消滅空間を $S\subset V_1,M\subset V_2$ とする。$p,z$ をつなぐ
$w(p,z)$ は非零である。零なら全ての $p,x_1,x_2,z$ が四文字閉経路となるからである。
これは $L_3$ 鉛筆の末尾からの縮約なので階数2である（非零な二項二次式の一次式倍は独立）。
$w(z)$ の $V_1$ を二項三次 $S^3 B$ と同定すると
$R_1\cap(V_1\otimes M)$ は二項二次式 $g$ に対し
$(s g,t g)$ を対応させる三次元行列空間である。$w(p,z)$ に対応する二次式を $g_0$ とすれば、
その第一支持が $S$ なので $S=g_0 B$ となる。従って $L$ と二項三次Veroneseの交叉は
$g_0$ の二根（重複込み）である。

さらに $R_1\cap(S\otimes M)$ はこの一本に等しい。
二次正方行列の二空間なら非零の可分解元があるからである。
従って $R_1$ は $R_1\cap(S\otimes V_2)$ と $R_1\cap(V_1\otimes M)$ の和五次元に、
両直線上の括弧を持ち上げる一本を足したものである。$q\in L$ で第一の空間は消え、
括弧の評価は $V_2/M$ で常に非零である。順側階数の低下は、第二空間の
$g_0$ 以外の二元を評価した $M$ 内の二行列式の零因子になる。
その行列は、$S^3B/(g_0 B)$ の双対から $g\mapsto (q(sg),q(tg))$ を取ったものなので、
その根は $q$ がVeronese上にある点ちょうどである（Hankel行列の階数1条件）。
両因子はともに長さ2で、二つの相異なる根なら各単根、一つなら二重根となる。
\end{proof}
